\documentclass{article}

\textwidth=6.5in
\textheight=8.5in
\voffset=-54pt
\oddsidemargin=0in
\evensidemargin=0in
\setlength{\parskip}{8pt}
\setlength{\parindent}{0pt}

\usepackage{amsmath, amssymb}
\usepackage{ifthen}

\usepackage[inline]{enumitem}
\setlist{topsep=0pt, parsep=8pt}

\usepackage[rgb,table]{xcolor}\convertcolorsUfalse
\usepackage{graphicx}

% change hyperref options
\PassOptionsToPackage{hyphens}{url}
\usepackage[bookmarks,colorlinks,linkcolor=black,citecolor=blue,pdfstartview=FitH,urlcolor=blue]{hyperref}
\hypersetup{pdfpagemode=UseNone}
\usepackage{bookmark}

\usepackage{graphicx}
\usepackage{multicol}
\usepackage{framed}
\usepackage{array}

\usepackage{icomma}

\usepackage{soul}

\usepackage{tikz}
\usetikzlibrary{
  matrix,%
  % enables the snake paths
  decorations.pathmorphing,%
  % enables bigger arrows
  decorations.markings,%
  decorations.pathreplacing,%
  decorations.text,%
  calc,%
  patterns,%
  shapes.geometric,%
  arrows,
  decorations.shapes,
  positioning,
  plotmarks,
  intersections
}
\tikzset{>=latex} % sets default arrow type in tikz
\tikzset{font=\normalsize}
\tikzset{mark size=1.5pt, mark options=thin}
\tikzset{pin distance=4pt,
  every pin edge/.style={<-, thin, shorten <= -2pt}}

\interfootnotelinepenalty=10000
\renewcommand{\thefootnote}{(\arabic{footnote})}

\usepackage{multirow, bigdelim}

% change to \usepackage[sols]{handout} to see solutions
\usepackage[sols]{handout}
\renewcommand{\workshopnote}{CA Notes.}    

\newcommand\iscurrentenumi[1]{%
  \ifthenelse{\equal{#1}{\theenumi}}{}{#1}}
\setenumerate[1]{ref=\#\arabic*}
\setenumerate[2]{ref=\protect\iscurrentenumi{\theenumi}(\alph*)}
\newcommand\enumiiprefix[2]{%
  \ifthenelse{{\equal{#1}{\theenumi}}\AND{\equal{#2}{(\alph{enumii})}}}{}{}%
  \ifthenelse{{\equal{#1}{\theenumi}}\AND{\NOT{\equal{#2}{(\alph{enumii})}}}}{#2}{}%
  \ifthenelse{\equal{#1}{\theenumi}}{}{#1#2}%
}
\setenumerate[3]{ref=\protect\enumiiprefix{\theenumi}{(\alph{enumii})}\roman*}

\makeatletter

% underline links               
\let\@href=\href
\renewcommand{\href}[2]{\@href{#1}{\ul{#2}}}

\makeatother

\definecolor{purple}{rgb}{0.5,0,1.0}

\usepackage{fancyhdr}
\lhead{\textcolor{white}{This content is protected and may not be shared, uploaded, or distributed.}}
\rhead{}
\rfoot{\vspace*{\fill}\footnotesize\textcopyright{} \emph{Harvard Math Ma}}
\renewcommand{\headrulewidth}{0pt}
\pagestyle{fancy}
\begin{document}

\htitle{Workshop 8}

\begin{workshop}
  Today's metacognitive focus in on checking your work / thinking about whether it's consistent.

  Give each student an index card, and have them answer the two questions on Slides 1 - 3 and then turn in their answers. The goal of having them turn in their answers is just so they commit to an answer. See the Google Slides for additional notes. 
\end{workshop}

\begin{enumerate}
\item 
  \begin{problem}[\stretch{0.5}]
    Your friend is working on the following problem.

    \begin{quote}
      A bacteria population is growing exponentially. At noon, the population has 400 bacteria. At 2 pm, it has 1200 bacteria. Write a formula giving the bacteria population $t$ hours after noon.
    \end{quote}

    Your friend has gotten the answer $400 \cdot 3^{2t}$. Does their answer work? How can you tell without solving the problem yourself?
  \end{problem} 

  \begin{solution}
      If I plug in $t=2$, I expect 1200 bacteria; however, plugging in $t=2$ yields $400(3^4)$ bacteria, which is much greater than 1200. The answer doesn't work; $B(t)=400(3)^{t/2}$.
  \end{solution}

\item 

  \begin{workshop}
    Here, the limits are inconsistent with the other information. 
  \end{workshop}

  \begin{problem}[\vfill\newpage]
    Another friend is trying to graph a continuous function $f(x)$. They've found that:
    \begin{itemize}
    \item $f(x)$ is increasing on $(-\infty, -2)$, decreasing on $(-2, 5)$, and increasing on $(5, \infty)$
    \item $f(x)$ is concave down on $(-\infty, 1)$ and concave up on $(1, \infty)$
    \item $\displaystyle \lim_{x \to -\infty} f(x) = \infty$ and $\displaystyle \lim_{x \to \infty} f(x) = 7$.
    \end{itemize}
    They ask for your help sketching the graph of $f$. What would you tell them?
  \end{problem}

  \begin{solution}
      \begin{itemize}
          \item It's not possible that $\displaystyle \lim_{x \to - \infty} f(x) = \infty$ because $f(x)$ increases on $(-\infty, 2)$. 
          \item It's not possible that $\displaystyle \lim_{x \to \infty}f(x) = 7$ because $f(x)$ increases and is concave up after $x=5$ ($\displaystyle \lim_{x \to \infty} f(x) = 7$ could be true if $f(x)$ increased and was concave down).
      \end{itemize}
  \end{solution}

\item 
  In each part, fill in the blanks with two consecutive integers, and explain your answer.
  \begin{enumerate}

  \item
    \begin{problem}[0.1in]
      If $3^A=7$, then $\cfitb{0.5in}{1} < A < \cfitb{0.5in}{2}$
    \end{problem}

    \begin{solution}
        $\log_3(3) < \log_3(7) < \log_3(9)$
    \end{solution}

  \item \label{estimate-ln2}
    \begin{problem}[0.1in]
      If $e^B=2$, then $\cfitb{0.5in}{0} < B < \cfitb{0.5in}{1}$
    \end{problem}

        \begin{solution}
        $\ln(1) < \ln(2) < \ln(e)$
    \end{solution}

  \item
    \begin{problem}[0.1in]
      If $e^C = \frac{1}{3}$, then $\displaystyle \cfitb{0.5in}{-2} < C < \cfitb{0.5in}{-1}$
    \end{problem}

            \begin{solution}
        $\ln\Big(\dfrac{1}{e^2}\Big) < \ln\Big(\dfrac{1}{3}\Big) < \ln\Big(\dfrac{1}{e}\Big)$
    \end{solution}

  \end{enumerate}



\item 

  \begin{workshop}
    We won't have talked about log rules yet, so the goal here is to use exponent rules. For example, in (c), $\dfrac{25}{\sqrt[3]{5}} = \dfrac{5^2}{5^{1/3}} = 5^{2 - 1/3} = 5^{5/3}$, so $\log_5 \left( \dfrac{25}{\sqrt[3]{5}} \right) = \dfrac{5}{3}$.
  \end{workshop}

  Simplify each of the following.

  \begin{enumerate}
  \item
    \begin{problem}[\vfill]
      Express the following value as 2 raised to a single exponent: $2^4 \cdot 2^5 $
    \end{problem}

    \begin{solution}
        $$2^4 \cdot 2^5 = 2^{4+5}=2^9$$
    \end{solution}

  \item
    \begin{problem}[\vfill]
      Express the following value as 3 raised to a single exponent: $ 9 \sqrt{3} $
    \end{problem}

    \begin{solution}
        $$9 \sqrt{3} = 3^2 \cdot 3^{1/2} = 3^{2+(1/2)} = 3^{5/2}$$
    \end{solution}

  \item
    \begin{problem}[\vfill]
      Express the following value as 5 raised to a single exponent: $\dfrac{25}{\sqrt[3]{5}}$
    \end{problem}

    \begin{solution}
        $$\dfrac{25}{\sqrt[3]{5}} = \dfrac{5^2}{5^{1/3}} = 5^{2-(1/3)} = 5^{5/3} $$
    \end{solution}

  \item
    \begin{problem}[\vfill\newpage]
      Express the following value as 7 raised to a single exponent: $ 49^5$
    \end{problem}

    \begin{solution}
        $$49^5 = (7^2)^5 = 7^{10}$$
    \end{solution}

  \end{enumerate}


\item 

  \begin{workshop}
    In this problems, students will need to have some idea of how big $\ln 2$ is. They figured this out already in \ref{estimate-ln2}, but it's worthwhile to let them think it through before you point them back there. 
  \end{workshop}

  Let $f(x) = 2x - e^x$.

  \begin{enumerate}
  \item

    \begin{problem}[\vfill]
      Where is $f$ increasing? Where is $f$ decreasing? 
    \end{problem}

  \item
    \begin{problem}[\vfill]
      Where is $f$ concave up? Where is $f$ concave down? 
    \end{problem}

  \item
    \begin{problem}[\vfill]
      Sketch a graph of $f$, labeling any local extrema with their exact coordinates. 
    \end{problem}

  \item

    \begin{problem}[0.25in]
      How many $x$-intercepts does $f$ have? How do you know?
    \end{problem}

    \begin{solution}
      Since the maximum value of $f$ is $2 \ln 2 - 2$ and we saw in \ref{estimate-ln2} that $\ln 2 < 1$, the maximum value of $f$ is negative. Therefore, $f$ has no $x$-intercepts.      
    \end{solution}

  \end{enumerate}

\end{enumerate}
\end{document}
